\chapter{Codetermination}
\label{sec:10}
A floor isn't just a list of forbidden acts — it's an allocation of authority, and there are four distinct powers in it: to \emph{write} the floor, \emph{amend} it, \emph{inspect} whether a deployed system still holds its commitment to it, and \emph{enforce} it by stopping the deployment. They usually shouldn't sit with the same institution, and the record shows what happens when they do, one power at a time.

Writing first. In February 2026 the terms under which Claude ran on the Pentagon's classified networks were Anthropic's to write, and the department demanded that they be rewritten, two prohibitions replaced by a license for all lawful use \autocite{defenseone2026ultimatum,volokh2026anthropiccontract}. The dispute ran between the two parties with a stake in the answer, and nobody the floor would protect had a say in what it would be.

Amendment next. Google's AI Principles, published in 2018 after the Maven protests, excluded weapons and surveillance \autocite{castelino2018google,conger2018maven,pichai2018principles} for nearly seven years until the pledge was removed in 2025 \autocite{google2025principles}. The demand became a \emph{company policy}: the objectors could hold the company to it in public, but the company that wrote it kept the power to amend it, so it lasted exactly as long as the company's willingness, and ending it had to be announced. That's the partial case: a transfer held at the institution's pleasure, and the 2025 announcement is what taking one back looks like. The pledge also bound only Google. The Maven work went to Palantir, and a refusal written into the standard license, the same for every vendor, would have gone with it.

Then inspection. In the Gebru and Mitchell dismissals, the company employed researchers to examine its own systems and dismissed them when the examination became costly \autocite{dickey2020gebru,dickey2021mitchell}.

Enforcement last. After the refusal, the President's order and the designation, Claude stayed in the targeting platform: in March Palantir said its products were still using it \autocite{cnbc2026karp}, and by September the department said about 90 percent of the work had moved off it \autocite{defensescoop2026michael}. Whatever its terms said, the vendor had no way to stop the deployment. Only the operator could, and it did so on its own timetable, which is where enforcement sits when nobody holds it apart.
